
\begin{example}
（2026海淀一模）已知椭圆$E:\frac{x^2}{a^2}+\frac{y^2}{b^2}=1(a > b > 0)$的离心率为$\frac{\sqrt{3}}{2}$，其左、右焦点和上顶点构成的三角形面积为$3\sqrt{3}$ \\
（1）求椭圆$E$的方程；\\
（2）设$O$为坐标原点，$M(x_0, y_0)$为第一象限内$E$上的动点，点$N$在直线$y=-2$上，且$OM\bot ON$，过点$A(0,3)$作$MN$的垂线交$OM$于$B$ ，求$\vert AB \vert$的值。
\end{example}

\begin{corollary*}
（1）$\frac{x^2}{12}+\frac{y^2}{3}=1$；（2） $\vert AB \vert =2\sqrt{3}$

\end{corollary*}


\begin{property}
（1）根据题意列出$a, b, c$的关系式，求出$a, b, c$ ，即可得到椭圆的标准方程 \\
（2）根据$OM\bot ON$结合点$N$在直线$y=-2$上求出点$N$的坐标；根据$MN\bot AB$得到$AB$的斜率，结合$A(0,3)$求出直线$AB$的方程；再求出直线$OM$的方程，$AB$与$OM$联立可得点$B$坐标；根据点与点的距离公式，利用$M$在椭圆上的约束条件计算$\vert AB \vert $。解题过程中需要补充讨论直线$AB$斜率不存在的特殊情形验证结论。

\begin{center}
\begin{tikzpicture}[
    node distance=4mm,
    box/.style={draw=customcolor, rounded corners=2pt,
        fill=customcolor!8, align=center, inner sep=3pt,
        text width=2.7cm, font=\small},
    result/.style={box, fill=customcolor!18},
    arrow/.style={->, semithick, customcolor!80!black}
]
    \node[box] (findN) {$OM\bot ON,\ y_N=-2$\\求点 $N$};
    \node[box, right=of findN] (lines)
        {$MN\bot AB,\ A(0,3)$\\$O,\ M(x_0,y_0)$\\求 $AB,\ OM$};
    \node[box, right=of lines] (findB) {$AB\cap OM$\\求点 $B$};
    \node[result, right=of findB] (distance)
        {距离公式 $+\ M\in E$\\$\lvert AB\rvert=2\sqrt{3}$};

    \draw[arrow] (findN) -- (lines);
    \draw[arrow] (lines) -- (findB);
    \draw[arrow] (findB) -- (distance);
\end{tikzpicture}
\end{center}

\end{property}

\begin{solution}
（1）由题意得$e= \frac{c}{a}= \frac{\sqrt{3}}{2}, \frac{1}{2}bc\times 4 = 3\sqrt{3}, a^2 = b^2 + c^2$
\\
所以$a = 2\sqrt{3}, b= \sqrt{3}, c= 3$\\
椭圆$E$的方程为$\frac{x^2}{12}+\frac{y^2}{3}=1$\\
（2）明显有$\frac{x_0^2}{12} + \frac{y_0^2}{3}=1(0< x_0 < 2\sqrt{3}, 0 < y_0 < \sqrt{3})$\\
由$OM \bot ON $可知$k_{OM}\cdot k_{ON} = -1$，即$k_{ON} = -\frac{x_0}{y_0}$\\
所以直线$ON: y = -\frac{x_0}{y_0}x$，令$y = -2$，得$N (\frac{2y_0}{x_0}, -2)$ \\
（i）当$\frac{2y_0}{x_0} \ne x_0$时，$k_{MN} = \frac{y_0 - (-2)}{x_0 - \frac{2y_0}{x_0}} = \frac{x_0(y_0 + 2)}{x_0^2 -2y_0} = \frac{x_0(y_0 + 2)}{12-4y_0^2 -2y_0} = \frac{x_0}{2(3 - 2y_0)}$
\begin{note}
    出现$x_0^2$，实现了代换降次
\end{note}
\noindent 所以$k_{AB} = \frac{2(2y_0 - 3)}{x_0}$\\
直线$AB$的方程为$y =\frac{2(2y_0 - 3)}{x_0}x + 3 $\\
直线$OM$的方程为$y = \frac{y_0}{x_0}x$\\
联立直线$AB$与$OM$得$B(\frac{x_0}{2-y_0}, \frac{y_0}{2-y_0})$\\
所以$\vert AB \vert = \sqrt{(\frac{x_0}{2-y_0})^2+(\frac{y_0}{2-y_0} - 3)^2} = \sqrt{\frac{x_0^2 + 16y_0^2 -48y_0 +36}{(2-y_0)^2}} = \sqrt{\frac{12-4y_0^2 + 16y_0^2 -48y_0 + 36}{(2-y_0)^2}} = 2\sqrt{3}$\\
（ii）当$\frac{2y_0}{x_0} = x_0$，此时直线$AB$斜率不存在，联立$\frac{x_0^2}{12} + \frac{y_0^2}{3}=1$得$x_0 = \sqrt{3}, y_0 = \frac{3}{2}$\\
直线$OM$的方程为$y = \frac{\sqrt{3}}{2}x$，直线$AB$的方程为$y = 3$
联立可得$B(2\sqrt{3}, 3)$，所以$\vert AB \vert = 2\sqrt{3}$\\
综上：$\vert AB \vert = 2\sqrt{3}$
\end{solution}



\begin{example}
（2025海淀一模）
已知椭圆$W:\frac{x^2}{a^2}+\frac{y^2}{b^2}=1(a > b > 0)$，$A, B$分别是$W$的左右顶点，$C$是$W$的上顶点，$\Delta ABC$的面积为2，且$\vert AC \vert = \sqrt{5}$ \\
（1）求椭圆$W$的方程及长轴长；\\
（2）已知点$M(2,1)$，点$P$在直线$AC$上，设直线$PM$与$x$轴交于点$E$，直线$BP$与直线$EC$交于点$Q$，判断点$Q$是否在椭圆$W$上，并说明理由。

\end{example}


\begin{corollary*}
（1）$\frac{x^2}{4} + y^2 = 1$，长轴长为4；（2）点$Q$在椭圆$W$上

\end{corollary*}
\begin{property}
（1）根据题意列出$a, b, c$的关系式，求出$a, b, c$ ，即可得到椭圆的标准方程\\（2）设出点$P$的坐标，求直线$PM$的方程并得点$E$坐标；求直线$BP$和$EC$的方程联立得点$Q$坐标，将其代入椭圆方程可判断点$Q$是否在椭圆上


\end{property}

\begin{solution}
（1）由题意得$\frac{1}{2} \cdot 2a \cdot b = 2, \sqrt{a^2 + b^2} = \sqrt{5}$\\
所以$a = 2, b = 1$\\
椭圆的方程为$\frac{x^2}{4} + y^2 = 1$，长轴长为$2a = 4$\\
（2）由（1）可得$A(-2, 0), B(2, 0), C(0, 1)$，直线$AC$的方程为$y = \frac{1}{2}x + 1$\\
设$P(2t, t+ 1)$
\begin{note}
    点$P$在直线上，所以可以只用一个未知数
\end{note}
（i）当$t \ne 1$且$t \ne 0$时，直线$PM$的方程为$y - 1 = \frac{t}{2t-2}(x - 2)$，令$y = 0$，得$E(\frac{2}{t}, 0)$ \\
直线$BP$的方程为$y = \frac{t+1}{2t - 2}(x - 2)$ \\
直线$EC$的方程为$y = -\frac{t}{2}x + 1$ \\
联立直线$BP$与直线$EC$可得$Q(\frac{4t}{t^2 + 1}, \frac{t^2 - 1}{t^2 + 1})$ \\
又$\frac{(\frac{4t}{t^2+1})^2}{4} + (\frac{t^2 - 1}{t^2 + 1})^2 = \frac{4t^2 + t^4 -2t^2 + 1}{(t^2 + 1)^2 } = \frac{(t^2 + 1)^2}{(t^2 + 1)^2} = 1$，即点$Q$在椭圆上\\
（ii）当$t = 1$时，此时直线$PM$的斜率不存在，点$Q$与点$B$重合，此时点$Q$在椭圆上 \\
（iii）当$t = 0$时，此时$PM$与$x$轴无交点，不符合题意 \\
综上：点$Q$在椭圆上
\end{solution}



\begin{example}
（2024海淀一模）已知椭圆$G: x^2 + my^2 =m$的离心率为$\frac{\sqrt{2}}{2}$，$A_1, A_2$分别是$G$的左、右顶点，$F$是$G$的右焦点\\
（1）求$m$的值及点$F$的坐标；\\
（2）设$P$是椭圆$G$上异于顶点的动点，点$Q$在直线$x = 2$上，且$PF \bot FQ$，直线$PQ$与$x$轴交于点$M$，比较$\vert MP \vert ^2$与$\vert MA_1\vert \cdot \vert MA_2 \vert$的大小

\end{example}

\begin{corollary*}
（1）$m = 2, F(1, 0)$；（2）$\vert MP \vert ^2 < \vert MA_1 \vert \cdot \vert MA_2 \vert$
    
\end{corollary*}

\begin{property}
（1）将椭圆方程转换为标准方程，根据离心率求出$m$,进一步确定$a, b, c$，可得焦点坐标 \\
（2）设点$P$坐标，根据$PF \bot FQ$结合点$Q$在直线$x = 2$上，求出点$Q$的坐标；求出直线$PQ$的方程并得点$M$坐标；根据点到点距离公式$\vert MP \vert ^2$与$\vert MA_1\vert \cdot \vert MA_2 \vert$，然后作差比较大小
    
\end{property}

\begin{solution}
（1）由$G : x^2 + my^2 = m$得$G: \frac{x^2}{m} + y^2 = 1$ \\
则有$e = \frac{c}{a} = \frac{\sqrt{m - 1}}{\sqrt{m}} = \frac{\sqrt{2}}{2}$，得$m = 2$ \\
故$G: \frac{x^2}{2} + y^2 = 1$，则$c = 1$，即$F(1, 0)$ \\
（2）由（1）可知$A_1(-\sqrt{2}, 0), A_2(\sqrt{2}, 0)$ \\
设$P(x_0, y_0)$，其中$x_0 \ne 0, y_0 \ne 0, -\sqrt{2} < x_0 < \sqrt{2}$，$Q(2, y_Q)$ \\
由$PF \bot FQ$可知$\overrightarrow{PF} \cdot \overrightarrow{FQ} = 0$，则$(x_0 - 1) \cdot (2 - 1) + y_0 y_Q = 0$，解得$y_Q = \frac{1 - x_0}{y_0}$，即$Q\left(2, \frac{1 - x_0}{y_0}\right)$ \\
直线$PQ$的斜率
\[
k_{PQ} = \frac{\frac{1 - x_0}{y_0} - y_0}{2 - x_0} = \frac{1 - x_0 - y_0^2}{y_0(2 - x_0)} = -\frac{x_0}{2y_0}
\]
直线$PQ$的方程为$y - y_0 = -\frac{x_0}{2y_0}(x - x_0)$，令$y = 0$，得$M\left(\frac{2}{x_0}, 0\right)$ \\
所以$\vert MP \vert^2 = \left(x_0 - \frac{2}{x_0}\right)^2 + y_0^2 = \frac{x_0^2}{2} + \frac{4}{x_0^2} - 3$，$\vert MA_1 \vert \cdot \vert MA_2 \vert = \left\vert \frac{2}{x_0} - \sqrt{2} \right\vert \cdot \left\vert \frac{2}{x_0} + \sqrt{2} \right\vert = \frac{4}{x_0^2} - 2$ \\
则$\vert MP \vert^2 - \vert MA_1 \vert \cdot \vert MA_2 \vert = \frac{x_0^2}{2} - 1$ \\
因为$-\sqrt{2} < x_0 < \sqrt{2}$，则$\frac{x_0^2}{2} - 1 < 0$ \\
所以$\vert MP \vert^2 - \vert MA_1 \vert \cdot \vert MA_2 \vert < 0$，即$\vert MP \vert^2 < \vert MA_1 \vert \cdot \vert MA_2 \vert$
    
\end{solution}


\begin{example}
（2024西城二模）已知椭圆$E:\frac{x^2}{a^2}+\frac{y^2}{b^2}=1(a > b > 0)$的一个顶点为$D(2, 0)$，焦距为$2\sqrt{2}$\\
（1）求椭圆$E$的方程；\\
（2）设点$P$是第一象限内椭圆$E$上一点，过$P$作$y$轴的垂线，垂足为$Q$，点$P$关于原点的对称点为$A$，直线$DQ$与椭圆的另一个交点为$B$，直线$DP$与$y$轴的交点为$C$，求证：$A, B, C$三点共线。
\end{example}

\begin{corollary*}
（1）$\frac{x^2}{4}+\frac{y^2}{2}=1$；（2）证明见详解

\end{corollary*}

\begin{property}
（1）由顶点坐标和焦距求出$a, c$，再由$b^2 = a^2 - c^2$得到椭圆方程 \\
（2）设点$P$的坐标，根据题意求出点$Q, A$的坐标；求出直线$DQ, DP$的方程，分别得到点$B, C$的坐标；再比较直线$AB$与$AC$的斜率，证明三点共线

\end{property}

\begin{solution}
（1）由题意得$a = 2, 2c = 2\sqrt{2}$，所以$c = \sqrt{2}$\\
则$b^2 = a^2 - c^2 = 4 - 2 = 2$，所以$b = \sqrt{2}$\\
椭圆$E$的方程为$\frac{x^2}{4}+\frac{y^2}{2}=1$\\
（2）设$P(x_0, y_0)$，则$x_0^2 + 2y_0^2 = 4(0 < x_0 < 2, 0 < y_0 < \sqrt{2})$，$Q(0, y_0)$，$A(-x_0, -y_0)$\\
直线$DQ$的方程为$y = -\frac{y_0}{2}(x - 2)$\\
设$B(x_B, y_B)$，联立直线$DQ$与椭圆$E$
\[
\begin{cases}
y = -\frac{y_0}{2}(x - 2) \\
\frac{x^2}{4} + \frac{y^2}{2} = 1
\end{cases}
\]
得$\left(1 + \frac{y_0^2}{2}\right)x^2 - 2y_0^2x + 2y_0^2 - 4 = 0$\\
则有$2 + x_B = \frac{2y_0^2}{1 + \frac{y_0^2}{2}}$，即$x_B = \frac{2(y_0^2 - 2)}{y_0^2 + 2}$\\
所以$y_B = -\frac{y_0}{2}(x_B - 2) = -\frac{y_0}{2}\left(\frac{2(y_0^2 - 2)}{y_0^2 + 2} - 2\right) = \frac{4y_0}{y_0^2 + 2}$，所以$B\left(\frac{2(y_0^2 - 2)}{y_0^2 + 2}, \frac{4y_0}{y_0^2 + 2}\right)$\\
直线$DP$的方程为$y = \frac{y_0}{x_0 - 2}(x - 2)$，令$x = 0$，得$C\left(0, -\frac{2y_0}{x_0 - 2}\right)$ \\
由$x_0^2 + 2y_0^2 = 4$得$y_0^2 = \frac{4 - x_0^2}{2}$\\
所以$x_B = \frac{2(y_0^2 - 2)}{y_0^2 + 2} = -\frac{2x_0^2}{8 - x_0^2}$，$y_B = \frac{4y_0}{y_0^2 + 2} = \frac{8y_0}{8 - x_0^2}$\\
则$k_{AB} = \frac{y_B - (-y_0)}{x_B - (-x_0)} = \frac{\frac{8y_0}{8 - x_0^2}+y_0}{-\frac{2x_0^2}{8 - x_0^2}+x_0} = \frac{y_0(16 - x_0^2)}{x_0(8 - x_0^2 - 2x_0)} = \frac{y_0(x_0 - 4)}{x_0(x_0 - 2)}$\\
$k_{AC} = \frac{-\frac{2y_0}{x_0 - 2} - (-y_0)}{0 - (-x_0)} = \frac{y_0(x_0 - 4)}{x_0(x_0 - 2)}$\\
所以$k_{AB} = k_{AC}$，即$A, B, C$三点共线

\end{solution}



\begin{example}
（2025东城二模）已知椭圆 $E:\frac{x^2}{a^2}+\frac{y^2}{b^2}=1\ (a>b>0)$ 的一个顶点为 $A(0,1)$，且过点 $\left(1,\frac{\sqrt{3}}{2}\right)$\\
(1) 求椭圆 $E$ 的方程及焦距；\\
(2) 过点 $\left(0,2\right)$ 的直线与椭圆 $E$ 交于不同的两点 $B,C$。直线 $AB,AC$ 的斜率分别记为 $k_1$ 与 $k_2$，当 $k_1+k_2=\frac{2}{3}$ 时，求 $\triangle ABC$ 的面积
\end{example}


\begin{corollary*}
（1）$\frac{x^2}{4}+y^2=1$，焦距为 $2\sqrt{3}$；（2）$\triangle ABC$ 的面积为 $\frac{2}{5}$
\end{corollary*}
\begin{property}
（1）根据题意列出$a, b, c$的关系式，求出$a, b, c$ ，即可得到椭圆的标准方程 \\
（2）点斜式设出$BC$的方程以及点$B, C$的坐标$B(x_1, y_1), C(x_2, y_2)$，将直线与椭圆联立得到韦达定理；用$x_1, x_2, y_1, y_2$表示直线$AB, AC$的斜率和；根据直线解析式把$y_1,y_2$统一替换为含$x_1,x_2$的式子，最后代入韦达定理求解得到斜率；根据弦长公式求出$\vert BC \vert$，再根据点到点距离公式求出$\vert BC \vert$边上的高，最后利用三角形面积公式计算$\triangle ABC$的面积。同时直线$BC$斜率不存在的情况需要单独讨论。

\end{property}

\begin{solution}
（1）由题意得：$b = 1, \frac{4}{a^2} + \frac{1}{b^2} = 1, a^2 = b^2 + c^2$\\
所以$a = 2, c = \sqrt{3}, b = 1$\\
椭圆的方程为$\frac{x^2}{4}+y^2=1$，焦距为$2c = 2\sqrt{3}$\\
（2）设直线$BC$的方程为$y = kx + 2$，$B(x_1, y_1), C(x_2, y_2)$\\
联立直线$BC$与椭圆$E$得
\[
\begin{cases}
y = kx + 2 \\
\frac{x^2}{4} + y^2 = 1
\end{cases}
\]
得$\left(1 + 4k^2\right)x^2 + 16kx + 12 = 0$，其中$\Delta = (16k)^2 - 4(1 + 4k^2)(12) = 256k^2 - 48 - 192k^2 = 64k^2 - 48 > 0$有$k^2 > \frac{3}{4}$\\
由韦达定理得$x_1 + x_2 = -\frac{16k}{1 + 4k^2}, x_1x_2 = \frac{12}{1 + 4k^2}$\\
所以$k_1 + k_2 = \frac{y_1 - 1}{x_1} + \frac{y_2 - 1}{x_2} = \frac{kx_1 + 1}{x_1} + \frac{kx_2 + 1}{x_2} = 2k + \frac{1}{x_1} + \frac{1}{x_2} = 2k + \frac{x_1 + x_2}{x_1x_2}$\\    
代入韦达定理可得$k_1 + k_2 = 2k - \frac{16k}{12} = \frac{2}{3}$，解得$k = 1$\\
此时直线$BC$的方程为$y = x + 2$，$x_1 + x_2 = -\frac{16}{5}, x_1x_2 = \frac{12}{5}$\\
所以$\vert BC \vert =  \sqrt{1 + k^2} \cdot \sqrt{(x_1 + x_2)^2 - 4x_1x_2} = \sqrt{2} \cdot \sqrt{\frac{256}{25} - \frac{48}{5}} = \frac{4\sqrt{2}}{5}$\\
又点$A(0,1)$到直线$BC$的距离为$d = \frac{|k \cdot 0 - 1 + 2|}{\sqrt{1 + k^2}} = \frac{1}{\sqrt{2}}$\\
所以$\triangle ABC$的面积为$S = \frac{1}{2} \cdot \vert BC \vert \cdot d = \frac{1}{2} \cdot \frac{4\sqrt{2}}{5} \cdot \frac{1}{\sqrt{2}} = \frac{2}{5}$\\
当直线$BC$斜率不存在时，即$k$不存在，此时直线$BC$为$x = 0$，与椭圆$E$的交点为$(0,1)$和$(0,-1)$，此时直线$AB$与直线$AC$的斜率之和不存在，不符合题意\\
综上：$\triangle ABC$的面积为$\frac{2}{5}$

\end{solution}



\begin{example}
（2026丰台一模）已知椭圆 $E:\frac{x^2}{a^2}+\frac{y^2}{b^2}=1\ (a>b>0)$，椭圆 $E$ 上一点 $A(2,1)$ 到两焦点的距离之和为 $2\sqrt{6}$\\
(1) 求椭圆 $E$ 的方程和离心率；\\
(2) 过点 $P(2,3)$ 的直线 $l$ 与椭圆 $E$ 交于不同的两点 $B,C$。若直线 $AB$ 与直线 $AC$ 的斜率之和为 $0$，求直线 $l$ 的方程
\end{example}

\begin{corollary*}
（1）$\frac{x^2}{6}+\frac{y^2}{3}=1$，离心率为 $\frac{\sqrt{2}}{2}$；（2）直线 $l$ 的方程为 $y = x + 1 $
\end{corollary*}
\begin{property}
（1）根据椭圆的定义得到$a$，点$A$在椭圆上，代入椭圆方程求出$b$，再由$c^2 = a^2 - b^2$得到$c$，最后计算离心率$e = \frac{c}{a}$ \\
（2）设直线$l$的方程为$x - 2 = n(y - 3)$，与椭圆方程联立得到韦达定理；用$x_1, x_2, y_1, y_2$表示直线$AB, AC$的斜率和；根据直线解析式把$x_1, x_2$统一替换为含$y_1, y_2$的式子，最后代入韦达定理求解得到斜率即得直线方程。同时直线$l$斜率为0的情况需要单独讨论。

\end{property}
\begin{solution}
（1）由题意可得$2a = 2\sqrt{6}, \frac{4}{a^2} + \frac{1}{b^2} = 1$\\
所以$a = \sqrt{6}, b = \sqrt{3}, c = \sqrt{3}$\\
椭圆的方程为$\frac{x^2}{6}+\frac{y^2}{3}=1$，离心率为$e = \frac{c}{a} = \frac{\sqrt{3}}{\sqrt{6}} = \frac{\sqrt{2}}{2}$\\
（2）设直线$l$的方程为$x - 2 = n(y - 3)$，点$B(x_1, y_1), C(x_2, y_2)$，直线与椭圆联立得
\[
\begin{cases}
x-2 = n(y - 3) \\
\frac{x^2}{6}+\frac{y^2}{3}=1
\end{cases}
\]
得$(n^2 + 2)y^2 +2n(2 - 3n) y +  9n^2 -12n -2 = 0$，其中$\Delta = (2n(2 - 3n))^2 - 4(n^2 + 2)(9n^2 - 12n - 2) = 4n^2(2 - 3n)^2 - 4(n^2 + 2)(9n^2 - 12n - 2) = -16(3n^2 -6n -1) > 0$\\
由韦达定理得$y_1 + y_2 = -\frac{2n(2 - 3n)}{n^2 + 2}, y_1y_2 = \frac{9n^2 -12n -2}{n^2 + 2}$\\
所以$k_1 + k_2 = \frac{y_1 - 1}{x_1 - 2} + \frac{y_2 - 1}{x_2 - 2} = \frac{y_1 - 1}{n(y_1 - 3)} + \frac{y_2 - 1}{n(y_2 - 3)} =0$ \\
即$\frac{y_1 - 1}{y_1 - 3} + \frac{y_2 - 1}{y_2 - 3} = 0$\\
整理可得$y_1y_2 - 2(y_1 + y_2) + 3 = 0$\\
代入韦达定理可得$$\frac{9n^2 -12n -2}{n^2 + 2} - 2\left(-\frac{2n(2 - 3n)}{n^2 + 2}\right) + 3 = 0$$
解得$n = 1$，满足$\Delta > 0$\\
所以直线$l$的方程为$x - 2 = y - 3$，即$y = x + 1$\\
当直线$l$斜率为0时，此时直线与椭圆无交点，不符合题意\\
综上：直线$l$的方程为$y = x + 1$
\end{solution}

\begin{example}

（2026朝阳一模）已知椭圆 $E:\frac{x^2}{a^2}+\frac{y^2}{b^2}=1\ (a>b>0)$ 的离心率为 $\frac{\sqrt{3}}{2}$，$A,B$ 分别为椭圆 $E$ 的上、下顶点，且 $B(0,-1)$\\
(1) 求椭圆 $E$ 的方程；\\
(2) 过点 $D(-1,-1)$ 的斜率存在且不为 $1$ 的直线 $l$ 与椭圆 $E$ 交于不同的两点 $M,N$（均不与点 $A$ 重合），点 $P$ 与点 $M$ 关于原点 $O$ 对称，直线 $BP$ 与直线 $OD$ 交于点 $Q$。求证：直线 $NQ$ 经过点 $A$
\end{example}

\begin{corollary*}
    （1）$\frac{x^2}{4}+y^2=1$；（2）证明见详解
\end{corollary*}
\begin{property}
（1）根据题意列出$a, b, c$的关系式，求出$a, b, c$ ，即可得到椭圆的标准方程 \\
（2）设直线$l$的方程为$y+1 = k(x+1)$，$M(x_1, y_1), N(x_2, y_2)$，将直线与椭圆联立得到韦达定理；根据对称关系求点$P$的坐标，再联立直线$BP$与直线$OD$的求出交点$Q$的坐标；最后利用三点共线知识证明直线$NQ$经过点$A$
\end{property}
\begin{solution}
（1）由题意得$b = 1, e = \frac{c}{a} = \frac{\sqrt{3}}{2} , a^2 = b^2 + c^2$，所以$a = 2, b = 1, c = \sqrt{3}$\\
椭圆的方程为$\frac{x^2}{4}+y^2=1$\\
（2）设直线$l$的方程为$y+1 = k(x+1)$，$M(x_1, y_1), N(x_2, y_2)$，将直线与椭圆联立得
\[\begin{cases}
y+1 = k(x+1) \\
\frac{x^2}{4}+y^2=1
\end{cases}\]
得$(4k^2 + 1)x^2 + 8k(k - 1)x + 4(k^2 -2k) = 0$，其中$\Delta = 64k^2(k-1)^2 -16(4k^2 + 1)((k^2 -2k)) = 16(3k^2 +2k) > 0$\\
由韦达定理得$x_1 + x_2 = -\frac{8k(k - 1)}{4k^2 + 1}, x_1x_2 = \frac{4(k^2 -2k)}{4k^2 + 1}$\\
点$P(-x_1, -y_1)$，直线$BP$的方程为$y  = \frac{-y_1 + 1}{-x_1}(x - 0) -1 = \frac{y_1 - 1}{x_1}x - 1$，直线$OD$的方程为$y = x$\\
联立直线$BP$与直线$OD$得$\frac{y_1 - 1}{x_1}x - 1 = x$，解得$x_Q = \frac{x_1}{y_1}$，所以$Q\left(\frac{x_1}{y_1 - x_1 - 1}, \frac{x_1}{y_1 - x_1 - 1}\right)$\\
所以$k_{AN}  - k_{AQ}= \frac{y_2 -1}{x_2} - \frac{\frac{x_1}{y_1 -x_1 - 1} - 1}{\frac{x_1}{y_1 -x_1 - 1} - 0} = \frac{(2k-2)x_1x_2 + (k-2)(x_1 + x_2)}{x_1x_2}$\\
而$(2k-2)x_1x_2 + (k-2)(x_1 + x_2) = (2k-2) \cdot \frac{4(k^2 -2k)}{4k^2 + 1} + (k-2) \cdot \left(-\frac{8k(k - 1)}{4k^2 + 1}\right) = 0$，即$k_{AN} = k_{AQ}$，故直线$NQ$经过点$A$，证毕

\end{solution}



\begin{example}
（2024朝阳期末）已知椭圆 $E:\frac{x^{2}}{a^{2}}+\frac{y^{2}}{b^{2}}=1\ (a>b>0)$ 的左顶点为 $A$，上顶点为 $B$，原点 $O$ 到直线 $AB$ 的距离为 $\frac{2\sqrt{5}}{5}$，$\triangle AOB$ 的面积为 $1$.\\
(1) 求椭圆 $E$ 的方程；\\
(2) 过点 $P(-2,1)$ 的直线 $l$ 与椭圆 $E$ 交于不同的两点 $C,D$，过点 $C$ 作 $x$ 轴的垂线分别与直线 $AD,AB$ 交于点 $M,N$. 判断点 $N$ 是否为线段 $CM$ 的中点，并说明理由.
\end{example}
